\section{Valor esperado de funciones de $X$ y $Y$}

\begin{teoria}{Definici\'on y propiedades}
\[
\E[g(X,Y)]=\sum_x\sum_y g(x,y)\,\fxy(x,y)\qquad\text{o}\qquad\E[g(X,Y)]=\iint g(x,y)\,\fxy(x,y)\,dx\,dy
\]
\guillemotleft Multiplica lo que quieres promediar por la conjunta y suma (o integra) sobre todo el soporte.\guillemotright
\begin{itemize}
\item $\E(c)=c$ \qquad $\E(aX+bY+c)=a\E(X)+b\E(Y)+c$ \quad(\textbf{siempre}, sean o no independientes).
\item S\'olo si son independientes: $\E(XY)=\E(X)\,\E(Y)$.
\end{itemize}
Casos que vas a usar: $g=x$ da $\E(X)$; $g=y$ da $\E(Y)$; $g=xy$ da $\E(XY)$; $g=x^2$ da $\E(X^2)$.
\end{teoria}

\begin{receta}{traducir el enunciado a $g(X,Y)$}
\begin{center}
\begin{tabular}{ll}
\guillemotleft total\guillemotright, \guillemotleft entre los dos\guillemotright & $g=X+Y$\\
\guillemotleft lo que queda\guillemotright, \guillemotleft sobrante\guillemotright & $g=Y-X$ (lo que hab\'ia menos lo que sali\'o)\\
\guillemotleft diferencia entre el m\'as caro y el m\'as barato\guillemotright & $g=|X-Y|$\\
\guillemotleft proporci\'on de una proporci\'on\guillemotright & $g=XY$\\
\guillemotleft costo $aX+b$\guillemotright & $g=aX+b$
\end{tabular}
\end{center}
\end{receta}

\begin{ejemplo}{Ejemplo resuelto 15 --- Queroseno \origen{Serie 3, problema 5}}
$Y$ = cantidad al inicio, $X$ = cantidad vendida, $\fxy=2$ en $0\le x\le y\le1$. Lo que queda es $Y-X$.
\[
\E(Y-X)=\int_0^1\!\!\int_0^y2(y-x)\,dx\,dy=\int_0^1\big[2yx-x^2\big]_0^{y}dy=\int_0^1y^2\,dy=\resp{\tfrac13}
\]
Quedan en promedio $\tfrac13$ de mil litros, unos 333 litros.
\end{ejemplo}

\begin{ejemplo}{Ejemplo resuelto 16 --- Impurezas \origen{Apuntes, Ej.~3.9}}
$\fxy=2-2x$ en el cuadrado unitario. La proporci\'on de impureza tipo I en la muestra es $XY$:
\[
\E(XY)=\int_0^1\!\!\int_0^1xy(2-2x)\,dy\,dx=\int_0^1x(1-x)\,dx=\tfrac12-\tfrac13=\resp{\tfrac16}
\]
\end{ejemplo}

\begin{ejemplo}{Ejemplo resuelto 17 --- Tabla independiente \origen{Serie 3, problema 4}}
\begin{center}
\begin{tabular}{c|cc|c}
$y\backslash x$ & 2 & 4 & $f_Y$\\ \hline
1 & 0.10 & 0.15 & 0.25\\
3 & 0.20 & 0.30 & 0.50\\
5 & 0.10 & 0.15 & 0.25\\ \hline
$f_X$ & 0.40 & 0.60 & 1
\end{tabular}
\end{center}
$\E(X)=2(0.4)+4(0.6)=3.2$,\quad $\E(Y)=1(0.25)+3(0.5)+5(0.25)=3$.

\textbf{a)} $\E(2X-3Y)=2(3.2)-3(3)=\resp{-2.6}$.

\textbf{b)} Como son independientes, $\E(XY)=\E(X)\E(Y)=3.2\cdot3=\resp{9.6}$.
(Comprobaci\'on directa: $2\cdot1\cdot0.10+4\cdot1\cdot0.15+2\cdot3\cdot0.20+4\cdot3\cdot0.30+2\cdot5\cdot0.10+4\cdot5\cdot0.15=9.6$.)
\end{ejemplo}

\begin{mecaniza}
\textbf{1.} (Autoexamen 5) Tabla: $f(0,1)=0.1$, $f(0,2)=0.3$, $f(1,1)=0.4$, $f(1,2)=0.2$. $\E(XY)$.
\quad\textbf{R:} $1\cdot1\cdot0.4+1\cdot2\cdot0.2=0.8$ (inciso B).

\textbf{2.} (Autoexamen 6) Dos l\'ineas de ensamble; n\'umero esperado de piezas de las dos.
$f_X=(0.35,0.37,0.28)$, $f_Y=(0.35,0.35,0.30)$.
\quad\textbf{R:} $\E(X)+\E(Y)=0.93+0.95=1.88$ (inciso D).
\end{mecaniza}

\subsection{Valor esperado condicional y curva de regresi\'on}

\begin{teoria}{Curva de regresi\'on}
Es el promedio de una variable cuando ya conozco la otra, usando la \textbf{condicional}:
\[
\hat y=\mu_{Y\mid x}=\E[Y\mid X=x]=\int y\,f_{Y\mid x}(y\mid x)\,dy
\qquad
\hat x=\mu_{X\mid y}=\E[X\mid Y=y]=\int x\,f_{X\mid y}(x\mid y)\,dx
\]
(con sumas si es discreta). El resultado es una funci\'on de la variable conocida; su gr\'afica es la curva de regresi\'on.
Sirve para pronosticar. Si son independientes, la curva es una constante: $\E[Y\mid x]=\E(Y)$.
\end{teoria}

\begin{receta}{curva de regresi\'on de $Y$ dado $x$}
(1) marginal $f_X$;\quad (2) condicional $f_{Y\mid x}=\fxy/f_X$;\quad (3) $\int y\,f_{Y\mid x}\,dy$ con los l\'imites de $y$;\quad
(4) simplifica y di si es una recta.
\end{receta}

\begin{ejemplo}{Ejemplo resuelto 18 \origen{Apuntes, Ej.~3.10 (en blanco)}}
Con $\fxy=x^3+\tfrac38xy^2$ y las condicionales del ejemplo resuelto 11:
\begin{align*}
\hat y=\E[Y\mid x]&=\int_0^2y\,\frac{8x^2+3y^2}{8(2x^2+1)}\,dy=\frac{16x^2+12}{8(2x^2+1)}=\resp{\frac{4x^2+3}{2(2x^2+1)}},\quad0\le x\le1\\[3pt]
\hat x=\E[X\mid y]&=\int_0^1x\,\frac{2x(8x^2+3y^2)}{4+3y^2}\,dx=\frac{\frac{16}{5}+2y^2}{4+3y^2}=\resp{\frac{2(5y^2+8)}{5(3y^2+4)}},\quad0\le y\le2
\end{align*}
No son rectas.
\end{ejemplo}

\begin{ejemplo}{Ejemplo resuelto 19 --- Regresi\'on lineal \origen{molde del Ej.~3.8}}
$\fxy=\tfrac83x(x+y)$, $0\le y\le x\le1$, con $f_X=4x^3$.
\[
f_{Y\mid x}=\frac{\frac83x(x+y)}{4x^3}=\frac{2(x+y)}{3x^2},\ 0\le y\le x
\qquad
\E[Y\mid x]=\int_0^xy\,\frac{2(x+y)}{3x^2}\,dy=\frac{2}{3x^2}\Big(\frac{x^3}{2}+\frac{x^3}{3}\Big)=\resp{\tfrac59x}
\]
Es una recta que pasa por el origen.
\end{ejemplo}

\begin{mecaniza}
\textbf{1.} (Serie 6c) Calcio y colesterol, densidad uniforme. Curvas de regresi\'on.
\quad\textbf{R:} hay independencia, as\'i que $\hat x=\E(X)=9.5$ y $\hat y=\E(Y)=180$: rectas constantes (una vertical y una horizontal).

\textbf{2.} (Serie 5) Con $\fxy=2$, $0\le x\le y\le1$: $\E[X\mid y]$ y $\E[Y\mid x]$.
\quad\textbf{R:} $f_Y=2y$, $f_{X\mid y}=\tfrac1y$ en $0\le x\le y$, $\E[X\mid y]=\tfrac y2$.
$f_X=2(1-x)$, $f_{Y\mid x}=\tfrac1{1-x}$ en $x\le y\le1$, $\E[Y\mid x]=\tfrac{1+x}2$. Ambas son rectas.

\textbf{3.} (Discreto, tabla del Ej.~3.3) $\E[X\mid Y=1]$.
\quad\textbf{R:} $0(0.2)+1(0.4)+2(0.4)=1.2$.
\end{mecaniza}

\subsection{Generadora de momentos de una suma (poco probable, pero es un rengl\'on)}

Si $X_1,\dots,X_n$ son \textbf{independientes} y $Y=X_1+\dots+X_n$, entonces
$M_Y(\theta)=M_{X_1}(\theta)\cdots M_{X_n}(\theta)$.

Ej.~3.11: $M_X=\dfrac{2}{2-\theta}$, $M_Y=\dfrac{1}{1-\theta}$, $T=X+Y$ $\Rightarrow$
$M_T(\theta)=\resp{\dfrac{2}{(2-\theta)(1-\theta)}}$.

\section{Covariancia y correlaci\'on}\label{sec:cov}

\begin{teoria}{Qu\'e miden}
\[
\Cov(X,Y)=\sigma_{XY}=\E(XY)-\E(X)\,\E(Y)
\qquad
\rho=\frac{\Cov(X,Y)}{\sqrt{\Var(X)}\sqrt{\Var(Y)}}=\frac{\sigma_{XY}}{\sigma_X\sigma_Y}
\qquad
\rho^2=\frac{\sigma_{XY}^2}{\sigma_X^2\sigma_Y^2}
\]
\begin{itemize}
\item \textbf{Signo de la covariancia:} positiva si cuando una sube la otra tiende a subir; negativa si una sube y la otra baja.
\item $\rho$ es la covariancia sin unidades: $-1\le\rho\le1$. Mide qu\'e tan \textbf{lineal} es la relaci\'on.
$\rho=\pm1$ s\'olo si $Y=mX+b$ exactamente.
\item $\rho^2$ (coeficiente de determinaci\'on): qu\'e fracci\'on de la variaci\'on de $Y$ se explica conociendo $X$.
\item Independientes $\Rightarrow\Cov=0$ y $\rho=0$. Al rev\'es no vale: $\Cov=0$ no garantiza independencia.
\end{itemize}
\end{teoria}

\begin{receta}{covariancia y correlaci\'on (las 5 esperanzas)}
\begin{enumerate}
\item Calcula $\E(X)$, $\E(Y)$, $\E(X^2)$, $\E(Y^2)$, $\E(XY)$.
(Las cuatro primeras salen m\'as r\'apido con las marginales.)
\item $\Var(X)=\E(X^2)-[\E(X)]^2$,\quad $\Var(Y)=\E(Y^2)-[\E(Y)]^2$.
\item $\Cov=\E(XY)-\E(X)\E(Y)$.
\item $\rho=\Cov/(\sigma_X\sigma_Y)$. Verifica que quede entre $-1$ y $1$.
\end{enumerate}
\end{receta}

\begin{ejemplo}{Ejemplo resuelto 20 --- Transistores \origen{Serie 3, problema 10b}}
Con la tabla del ejemplo resuelto 5: $f_X=f_Y=(0.6,\,0.3,\,0.1)$ en los valores $1,2,3$.
\begin{align*}
\E(X)&=\E(Y)=1(0.6)+2(0.3)+3(0.1)=1.5\\
\E(X^2)&=\E(Y^2)=1(0.6)+4(0.3)+9(0.1)=2.7\\
\Var(X)&=\Var(Y)=2.7-(1.5)^2=0.45\\
\E(XY)&=1(0.3)+2(0.2)+3(0.1)+2(0.2)+4(0.1)+3(0.1)=2.1\\
\Cov(X,Y)&=2.1-(1.5)(1.5)=-0.15\qquad\qquad\rho=\frac{-0.15}{\sqrt{0.45}\sqrt{0.45}}=\resp{-\tfrac13}
\end{align*}
Tiene sentido que sea negativa: si el primer defectuoso tarda en salir, quedan menos piezas y el segundo sale pronto.
\end{ejemplo}

\begin{ejemplo}{Ejemplo resuelto 21 --- Neum\'aticos \origen{Serie 3, problema 9}}
$\fxy=K(x^2+y^2)$, $K=\tfrac{3}{380000}$, en $[20,30]\times[20,30]$, con $f_X(x)=\tfrac{3x^2}{38000}+\tfrac1{20}$
(y lo mismo para $Y$ por simetr\'ia).
{\small\begin{align*}
\E(X)&=\int_{20}^{30}x\Big(\tfrac{3x^2}{38000}+\tfrac1{20}\Big)dx=\tfrac{1925}{76}=25.3289=\E(Y)\\
\E(X^2)&=\int_{20}^{30}x^2\Big(\tfrac{3x^2}{38000}+\tfrac1{20}\Big)dx=\tfrac{37040}{57}=649.8246\ \Rightarrow\ \Var(X)=\Var(Y)=8.2690\\
\E(XY)&=K\int_{20}^{30}\!\!\int_{20}^{30}xy(x^2+y^2)\,dx\,dy=\tfrac{24375}{38}=641.4474
\end{align*}}
\textbf{a)} $\Cov(X,Y)=641.4474-(25.3289)^2=\resp{-0.1082}$ (la respuesta $-0.108$ de la serie).

\textbf{b)} $\rho=\dfrac{-0.1082}{8.2690}=\resp{-0.0131}$. Relaci\'on lineal pr\'acticamente nula.

\textbf{Cuidado:} aqu\'i se restan dos n\'umeros casi iguales. Trabaja con fracciones o con al menos 4 decimales; si redondeas
$\E(X)$ a 25.33 la covariancia sale mal.
\end{ejemplo}

\begin{ejemplo}{Ejemplo resuelto 22 \origen{Apuntes, Ej.~3.12 (en blanco), incisos c y d}}
$\fxy=4xye^{-(x^2+y^2)}$. Ya vimos que son independientes con $f_X=2xe^{-x^2}$.

\textbf{c)} Por independencia las esperanzas condicionales son constantes:
$\E[X\mid y]=\E(X)=\int_0^\infty2x^2e^{-x^2}dx=\resp{\tfrac{\sqrt\pi}{2}}$, y $\E[Y\mid x]=\resp{\tfrac{\sqrt\pi}{2}}$.

\textbf{d)} Independientes $\Rightarrow\Cov(X,Y)=0\Rightarrow\resp{\rho=0}$. No hace falta integrar.
\end{ejemplo}

\begin{mecaniza}
\textbf{1.} (Autoexamen 7) $\fxy=4xy$ en el cuadrado unitario. $\Cov(X,Y)$.
\quad\textbf{R:} se factoriza como $(2x)(2y)$ en un rect\'angulo: independientes, $\Cov=0$ (inciso C).

\textbf{2.} (Autoexamen 10) $f(0,1)=0.1$, $f(1,1)=0.3$, $f(0,2)=0.4$, $f(1,2)=0.2$. Halla $\rho$.
\quad\textbf{R:} $\E(X)=0.5$, $\E(Y)=1.6$, $\E(XY)=0.3+0.4=0.7$, $\Cov=-0.1$, $\Var(X)=0.25$, $\Var(Y)=0.24$,
$\rho=-0.1/\sqrt{0.06}=-0.408$ (inciso A).

\textbf{3.} Con $\fxy=24xy$ del tri\'angulo $x+y\le1$: $\Cov$ y $\rho$.
\quad\textbf{R:} $\E(X)=\E(Y)=\tfrac25$, $\E(X^2)=\tfrac15$, $\Var=\tfrac1{25}$, $\E(XY)=\tfrac2{15}$,
$\Cov=\tfrac2{15}-\tfrac4{25}=-\tfrac2{75}$, $\rho=-\tfrac23$.

\textbf{4.} Con $\fxy=x+y$ del Ej.~3.4: $\Cov$ y $\rho$.
\quad\textbf{R:} $\E(X)=\E(Y)=\tfrac7{12}$, $\E(XY)=\tfrac13$, $\Cov=-\tfrac1{144}$, $\Var=\tfrac{11}{144}$, $\rho=-\tfrac1{11}$.
\end{mecaniza}

\section{Combinaciones lineales}

\begin{teoria}{Las dos f\'ormulas}
Para $W=aX+bY+c$:
\[
\E(W)=a\,\E(X)+b\,\E(Y)+c
\qquad
\Var(W)=a^2\Var(X)+b^2\Var(Y)+2ab\,\Cov(X,Y)
\]
\begin{itemize}
\item La constante $c$ \textbf{no cambia} la variancia.
\item Los coeficientes entran \textbf{al cuadrado} en la variancia: un signo menos desaparece ah\'i, pero \textbf{no} en el t\'ermino $2ab\Cov$.
\item Si son independientes, $\Cov=0$ y queda $\Var(W)=a^2\Var(X)+b^2\Var(Y)$.
\item Si te dan $\rho$ en vez de la covariancia: $\Cov=\rho\,\sigma_X\sigma_Y$.
\end{itemize}
Con $n$ variables: $\Var\big(\sum a_iX_i\big)=\sum a_i^2\Var(X_i)+2\sum_{i<j}a_ia_j\Cov(X_i,X_j)$, o en forma matricial
$\mathbf a^{T}\Sigma\,\mathbf a$ con $\Sigma$ la matriz de variancia-covariancia.
\end{teoria}

\begin{ejemplo}{Ejemplo resuelto 23 \origen{Serie 3, problema 11}}
$X$, $Y$ independientes, $\E(X)=3$, $\E(X^2)=25$, $\E(Y)=10$, $\E(Y^2)=164$.

\textbf{a)} $\E(3X+Y-8)=9+10-8=\resp{11}$\qquad
\textbf{b)} $\E(2X-3Y+7)=6-30+7=\resp{-17}$

\textbf{c)} $\Var(X)=25-9=\resp{16}$\quad \textbf{d)} $\sigma_X=\resp{4}$\quad
\textbf{e)} $\Var(Y)=164-100=\resp{64}$\quad \textbf{f)} $\sigma_Y=\resp{8}$

\textbf{g)} $\Var(3X+Y-8)=9(16)+64=\resp{208}$\qquad
\textbf{h)} $\Var(2X-3Y+7)=4(16)+9(64)=\resp{640}$

\textbf{i)} $\E\big[\tfrac{X-3}{4}\big]=\tfrac{3-3}{4}=\resp{0}$,\quad $\Var\big[\tfrac{X-3}{4}\big]=\tfrac{1}{16}(16)=\resp{1}$

\textbf{j)} $\E\big[\tfrac{Y-10}{8}\big]=\resp{0}$,\quad $\Var\big[\tfrac{Y-10}{8}\big]=\tfrac1{64}(64)=\resp{1}$

\textbf{k)} S\'i: para cualquier variable, $Z=\dfrac{X-\mu}{\sigma}$ tiene
$\E(Z)=\dfrac{\mu-\mu}{\sigma}=0$ y $\Var(Z)=\dfrac{\sigma^2}{\sigma^2}=1$ (estandarizaci\'on).
\end{ejemplo}

\begin{ejemplo}{Ejemplo resuelto 24 --- Tres variables \origen{Apuntes, Ej.~3.13 (en blanco)}}
$\E(Y_1)=2$, $\E(Y_2)=-1$, $\E(Y_3)=4$; $\Var=4,\,6,\,8$; $\Cov(Y_1,Y_2)=1$, $\Cov(Y_1,Y_3)=-1$, $\Cov(Y_2,Y_3)=0$.
Para $W=3Y_1+4Y_2-6Y_3$:
\begin{align*}
\E(W)&=3(2)+4(-1)-6(4)=\resp{-22}\\
\Var(W)&=\underbrace{9(4)+16(6)+36(8)}_{420}+2\big[\underbrace{(3)(4)(1)}_{12}+\underbrace{(3)(-6)(-1)}_{18}+\underbrace{(4)(-6)(0)}_{0}\big]=420+60=\resp{480}
\end{align*}
Forma matricial (da lo mismo):
$\begin{bmatrix}3&4&-6\end{bmatrix}
\begin{bmatrix}4&1&-1\\1&6&0\\-1&0&8\end{bmatrix}
\begin{bmatrix}3\\4\\-6\end{bmatrix}=480$.
\end{ejemplo}

\begin{mecaniza}
\textbf{1.} (Autoexamen 8) $\sigma_X^2=3$, $\sigma_Y^2=4$, $\Cov=2$. $\Var(2X-Y)$.
\quad\textbf{R:} $4(3)+1(4)+2(2)(-1)(2)=12+4-8=8$ (inciso E).

\textbf{2.} (Autoexamen 9) $\mu_X=10$, $\sigma_X=15$, $\mu_Y=15$, $\sigma_Y=20$, $\rho=-0.1$; $R=0.4X+0.6Y$.
\quad\textbf{R:} $\E(R)=4+9=13$. $\Cov=-0.1(15)(20)=-30$.
$\Var(R)=0.16(225)+0.36(400)+2(0.4)(0.6)(-30)=36+144-14.4=165.6$, $\sigma_R=12.87$ (inciso D).
\end{mecaniza}

\begin{trampa}
\begin{itemize}
\item Te dan $\sigma$ (desviaci\'on) y la usas como variancia. En la f\'ormula va $\sigma^2$.
\item $\Var(2X-3Y)$ \textbf{no} es $4\Var(X)-9\Var(Y)$: es $4\Var(X)+9\Var(Y)-12\Cov(X,Y)$.
\item Sumar la constante a la variancia. $\Var(X+7)=\Var(X)$.
\item Usar $\E(XY)=\E(X)\E(Y)$ sin que el problema diga (o hayas probado) que son independientes.
\end{itemize}
\end{trampa}
